\documentclass[11pt]{article}
\usepackage[margin=1in]{geometry}
\usepackage{xcolor}
\usepackage{enumitem}
\usepackage{booktabs}
\usepackage{hyperref}
\usepackage{amsmath}
\usepackage{amssymb}

\definecolor{reviewercolor}{RGB}{0,0,180}
\definecolor{responsecolor}{RGB}{0,100,0}

\newcommand{\reviewer}[1]{\textcolor{reviewercolor}{\textbf{Reviewer Comment:} #1}}
\newcommand{\response}[1]{\textcolor{responsecolor}{\textbf{Our Response:} #1}}
\newcommand{\change}[1]{\textit{Change in manuscript:} #1}

\title{Response to Reviewers\\
\large Texture Resonance Retrieval for Retrieval-Grounded\\Editable Audio Effect Preset Selection\\
\large IEEE Transactions on Multimedia --- Major Revision}
\author{Anonymous Authors}
\date{}

\begin{document}
\sloppy
\maketitle

We sincerely thank the reviewers for their detailed and constructive feedback. Below we address each concern point by point, organized by the major (M1--M5) and minor (m1--m8) issues identified in the review. All referenced section, table, and figure numbers refer to the revised manuscript.

\section*{Major Issues}

\subsection*{M1: Limited Dataset Scale and Near-Duplicate Risk}

\reviewer{The benchmark comprises only 1,267 guitar presets (204 test queries), which is small by retrieval standards. After removing exact audio reuse, 200 near-duplicate pairs remain. This raises the concern that strong retrieval performance may be inflated by trivially close knowledge-base items.}

\response{We agree that dataset scale is a genuine limitation and have taken concrete steps to quantify its impact on the reported results:}

\begin{enumerate}[leftmargin=*]
\item \textbf{Near-duplicate sensitivity analysis} (new Protocol-A paragraph ``Near-duplicate sensitivity''): We progressively remove near-duplicates from the knowledge base at four L2 thresholds ($\tau \in \{0.005, 0.01, 0.02, 0.05\}$) and re-evaluate TRR. The normalized L2 degrades as the KB shrinks from 1,063 to 83 candidates, so the revised paper treats database density as a real boundary condition rather than as a solved issue. Full details appear in the supplementary sensitivity subsection.
\item \textbf{Near-duplicate detection infrastructure}: We provide a reusable \texttt{near\_dup\_filter.py} script and \texttt{near\_dup\_sensitivity.py} experiment, enabling reviewers to reproduce the analysis at arbitrary thresholds.
\item \textbf{Explicit scope acknowledgment}: We clearly state in the experimental setup, Discussion, and Limitations that this is a pilot benchmark on guitar effects, not a claim of universal scale.
\end{enumerate}

\change{Protocol-A near-duplicate sensitivity paragraph and table added; supplementary sensitivity analysis retained; Abstract revised to mention near-duplicate sensitivity.}

\subsection*{M2: Weak Baseline Coverage}

\reviewer{The paper compares TRR only against text-RAG, FeatureNN-RAG, Wav2Vec-RAG, and CLAP. Strong audio representation baselines such as PANNs and PaSST, and a direct regression baseline (MLP), are absent.}

\response{We have added three new baselines:}

\begin{enumerate}[leftmargin=*]
\item \textbf{PANNs (CNN14)}: Large-scale pretrained audio pattern recognition network [Kong et al., 2020]. We use the CNN14 variant with AudioSet-pretrained embeddings for retrieval.
\item \textbf{PaSST}: Efficient audio transformer with Patchout [Koutini et al., 2022]. Retrieval uses the AudioSet-pretrained embedding space.
\item \textbf{MLP-Regressor}: A two-layer MLP with 256 hidden units trained on mean-pooled Wav2Vec2 features to directly predict flattened parameter vectors, establishing a supervised regression boundary.
\end{enumerate}

All retrieval baselines are evaluated under Protocol-A with the same split and metrics. Results are reported in Table~3 of the main paper, while the MLP boundary is reported separately because it predicts dense parameter leaves rather than retrieving executable presets.

\change{Table~3 expanded with PANNs and PaSST retrieval rows; direct-regression boundary table added; baseline descriptions updated.}

\subsection*{M3: No Ablation Studies}

\reviewer{There are no ablation studies validating the core design choices of TRR: projection dimensionality, layer selection, and projection type (PCA vs.\ random).}

\response{We have added a controlled supplementary mechanism-ablation subsection with three diagnostic comparisons:}

\begin{enumerate}[leftmargin=*]
\item \textbf{Projection dimensionality and layer selection}: Supplementary Table~S2 reports controlled mechanism rows including $d{=}32$ and $d{=}64$ projected multi-layer variants, layer~5 Gram aggregation, and the full unprojected Gram boundary.
\item \textbf{Aggregation and normalization}: Supplementary Table~S2 compares same-backbone mean pooling, Gram aggregation, and a no-normalization variant. This isolates the contribution of second-order aggregation without claiming global architectural optimality.
\item \textbf{Interpretation guardrail}: The main paper now cites the ablation only as mechanism support for TRR as a retrieval prior, not as proof that the chosen layer set and projection dimension are globally optimal.
\end{enumerate}

\change{Supplementary mechanism-ablation subsection and Table~S2 added; Algorithm 1 updated to specify ``frozen random linear projection'' and correct dimension to 32; Abstract mentions mechanism ablations.}

\subsection*{M4: Unnormalized Metrics}

\reviewer{All L2 error metrics operate on raw parameter values with heterogeneous physical scales. A compressor threshold in dB and a delay time in ms have very different magnitudes, making raw L2 hard to interpret and potentially dominated by large-scale parameters.}

\response{We now report a \textbf{Normalized L2} (Norm.\,L2) metric alongside the original raw L2:}

\begin{enumerate}[leftmargin=*]
\item Each parameter leaf is min-max normalized to $[0,1]$ using the physical DSP parameter ranges defined in the plugin source code (\texttt{PluginAudioParameters.h}).
\item The ranges are documented in a shareable \texttt{param\_ranges.json} file for reproducibility.
\item The evaluation pipeline (\texttt{evaluate.py}) supports both modes: \texttt{Evaluator(normalize=False)} for backward compatibility and \texttt{Evaluator(normalize=True)} for the new normalized metrics.
\end{enumerate}

\change{Section 4.2 ``Normalized L2'' paragraph (new); Table~3 adds a Norm.\,L2 column; \texttt{param\_ranges.json} and updated \texttt{evaluate.py} provided.}

\subsection*{M5: Listening Test Methodological Weaknesses}

\reviewer{Trials 2--5 compare the system against a ``manual'' baseline with unknown provenance and time budget. There are also no demographic details about participant expertise.}

\response{We have addressed this in several ways:}

\begin{enumerate}[leftmargin=*]
\item \textbf{Condensed Trials 2--5}: The main text now reports aggregate trial-block summaries with explicit caveats about manual-baseline provenance, rather than treating the manual comparison as a fairness-critical optimization benchmark.
\item \textbf{Participant metadata boundary}: We report the available age and gender metadata and explicitly state that the logs do not preserve device metadata, expertise, or manual-baseline time budget.
\item \textbf{Local vLLM diagnostic}: Following the revision request to add evidence where feasible, we added a conservative local qwen3.5-9b vLLM diagnostic that compares text-side perceptual-recognizability scores with Trial-1 human ratings. The result is negative and is reported as a boundary condition rather than as supportive perceptual evidence.
\item \textbf{Interpretation guardrails}: We explicitly frame the listening test as exploratory perceptual context, not as validation of the objective metrics or proof of superiority over MusicGen.
\end{enumerate}

\change{Section 4.5 condensed and reframed; local vLLM diagnostic added; missing participant-expertise and manual-baseline provenance now stated as limitations.}

\section*{Minor Issues}

\subsection*{m1: Unclear Algorithm 1 Projection Description}

\reviewer{Algorithm 1 describes a ``linear projection $P: \mathbb{R}^{768} \to \mathbb{R}^{64}$'' but the code uses $d{=}32$.}

\response{Fixed. Algorithm 1 now correctly states ``frozen random linear projection $P: \mathbb{R}^{768} \to \mathbb{R}^{32}$'' with Xavier initialization and no gradient updates. All downstream dimension references (Gram matrix $32 \times 32$, embedding 1024-dimensional) are consistent.}

\change{Algorithm 1 corrected; prose in Section 3.2 updated.}

\subsection*{m2: Over-Hedging}

\reviewer{Phrases like ``currently evaluated,'' ``we therefore interpret,'' and ``on the current benchmark'' appear excessively, reducing readability.}

\response{We conducted a systematic pass to reduce redundant qualifiers while preserving scientific rigor. Key reductions include removing unnecessary scope modifiers, simplifying ``We therefore interpret X as...'' constructions, and keeping only the qualifiers needed to state the benchmark boundary honestly.}

\change{Throughout the manuscript; redundant hedging reduced while retaining necessary scope boundaries.}

\subsection*{m3: Section 3.3 (Preliminary Personalization Extension) is Unreferenced}

\reviewer{Section 3.3 describes a personalization system but no results support it, and it contributes to page bloat.}

\response{Deleted entirely. The Discussion now explicitly excludes personalization from validated claims.}

\change{Section 3.3 removed.}

\subsection*{m4: Table 1 (Parameter--Waveform Duality) Lacks Data}

\reviewer{Table 1 contains only a conceptual comparison with no empirical data and consumes valuable page budget.}

\response{Table 1 has been removed. The Parameter--Waveform Duality concept is retained as a brief inline paragraph in Section 2.}

\change{Table 1 removed; inline description preserved.}

\subsection*{m5: Missing Visualization}

\reviewer{No visualizations of the TRR embedding space or Gram matrix structure are provided.}

\response{We have added two figures:}
\begin{enumerate}[leftmargin=*]
\item \textbf{t-SNE diagnostic}: comparison of TRR vs.\ Wav2Vec2 mean-pooled embeddings, colored by style label.
\item \textbf{Gram-matrix diagnostic}: representative Gram-matrix examples showing the second-order correlation structure captured by TRR.
\end{enumerate}

\change{Two vector diagnostic figures added and referenced in Section 3.2.}

\subsection*{m6: BibTeX Errors}

\reviewer{Several BibTeX entries have wrong types (e.g., \texttt{@inproceedings} for journal articles like Neuron), incomplete author lists (``and others''), and missing metadata.}

\response{Targeted BibTeX cleanup performed for the entries most likely to affect reviewer trust: \texttt{LLM2Fx} is no longer represented as an arXiv ``proceedings'' item, its author list and arXiv DOI are expanded, \texttt{AudioLDM} is represented as an ICML 2023 paper, and \texttt{MusicLM}'s arXiv entry now uses the full arXiv author list rather than ``and others.'' We also corrected the previously noted journal-entry type issue for \texttt{mcdermott2011sound}.}

\change{\texttt{reference.bib}: high-risk bibliographic entries corrected; remaining venue metadata should be checked again before final camera-ready submission.}

\subsection*{m7: Figure 1 Caption is Vague}

\reviewer{The system diagram caption does not specify the DSP effect chain topology.}

\response{Updated to explicitly list the nine-module topology: Compressor $\to$ Driver $\to$ Screamer $\to$ Equaliser $\to$ Chorus $\to$ Flanger $\to$ Phaser $\to$ Delay $\to$ Reverb, with each module switchable on/off.}

\change{Figure 1 caption rewritten.}

\subsection*{m8: Missing Fusion Score Formula}

\reviewer{The fusion mechanism is described only verbally; no equation is provided.}

\response{Added explicit formula: $s_{\mathrm{fused}}(z_i)=\alpha s_{\mathrm{text}}(z_i)+(1-\alpha)s_{\mathrm{audio}}(z_i)$, where $\alpha$ is the text-side weight after score normalization. The manuscript now states that the adaptive profile is evaluated only as a robustness diagnostic.}

\change{Fusion formula added in the ``Multimodal Fusion as a Supplementary Diagnostic Extension'' subsection.}

\section*{Summary of Changes}

\begin{table}[h]
\centering
\small
\begin{tabular}{lll}
\toprule
\textbf{Issue} & \textbf{Type} & \textbf{Action Taken} \\
\midrule
M1 Dataset scale & Major & Near-dup sensitivity analysis + supplementary table \\
M2 Weak baselines & Major & PANNs, PaSST, MLP-Regressor added \\
M3 No ablations & Major & 3 ablation experiments + subsection \\
M4 Unnormalized metrics & Major & Norm.\,L2 column + param\_ranges.json \\
M5 Listening test issues & Major & Condensed and reframed as exploratory perceptual context \\
m1 Algorithm dim error & Minor & Algorithm 1 corrected to $d{=}32$ \\
m2 Over-hedging & Minor & 15+ instances reduced \\
m3 Unreferenced Sec 3.3 & Minor & Section deleted \\
m4 Table 1 no data & Minor & Table removed \\
m5 Missing visualization & Minor & t-SNE + Gram heatmap figures added \\
m6 BibTeX errors & Minor & High-risk entries corrected \\
m7 Vague Fig 1 caption & Minor & Caption rewritten with topology \\
m8 Missing fusion formula & Minor & Formula added \\
\bottomrule
\end{tabular}
\end{table}

We believe these revisions substantially strengthen the paper and address all reviewer concerns. We remain happy to provide further clarification or additional experiments if needed.

\end{document}
